% Lecture visuelle directe avec la vision du modèle ; aucun OCR.
% Source : PDF original, vue 019, côté gauche, page manuscrite 38.
% Statut : transcrit ; relecture indépendante à effectuer.
\chapter{Journée du 18 brumaire et signature du concordat}

Continuons, MM., à observer la conduite miséricordieuse de la Providence envers nous ; elle vient de manifester d'une manière éclatante sa protection sur l'Église universelle ; maintenant elle va tout préparer pour annoncer bientôt à l'Église de France des jours plus prospères.

C'était le gal Bonaparte qui, par les ordres du Directoire, avait fait la conquête de l'Italie ; ce personnage s'était servi de toute l'influence qu'avait pu lui donner la victoire pour seconder la tyrannie qui pesait alors sur la France. Mais depuis, l'ambition qui plus tard l'a rendu si célèbre sous le nom de l'empereur Napoléon commençait à paraître. Les directeurs, pour se délivrer de l'inquiétude que leur causait un homme qui, à la tête d'une armée victorieuse, aujourd'hui leur appui, demain pouvait devenir leur maître, imaginèrent de le rappeler de l'Italie sous le prétexte honorable de lui donner le commandement d'une expédition qu'ils dirigeaient contre l'Égypte. Dieu qui dispose de tout pour arriver à ses fins le permettait ainsi, car c'était le moment où l'armée révolutionnaire allait être chassée de l'Italie pour y laisser rétablir l'autorité du chef de l'Église, et pour que les décrets du Ciel pussent s'accomplir plus tard il fallait que Bonaparte ne se
% Source : PDF original, vue 019, côté droit, page manuscrite 39.
trouvât pas à la tête de l'armée vaincue. Le maître souverain le tenait en réserve jusqu'au moment où il devait être chargé d'une double mission, mission de miséricorde envers la France en y rétablissant la religion et l'ordre ; mission de justice envers les nations en visitant, à main armée, toutes leurs capitales. Nous ne le suivrons point en Égypte ; qu'il nous suffise de savoir qu'après quelques succès passagers, l'expédition échoua complètement. Bonaparte abandonne son armée et revient subitement à Paris ; il tente alors une audacieuse entreprise qui est couronnée du succès le plus complet. Favorisé par la majorité du Conseil des Anciens et par le concours de son frère Lucien qui se trouve pour lors président du Conseil des Cinq-Cents, il disperse ce qui s'appelaient les représentants du peuple, renverse le Directoire, brise la constitution de la République et la remplace par une constitution nouvelle qui, sous le nom de premier consul, met entre ses mains l'autorité suprême, en attendant le moment où plus tard il osera se revêtir de la pourpre impériale. C'est cette révolution qu'on a appelée le 18 brumaire. Ces événements se passèrent dans les premiers jours de novembre 1799. S'ils ne rendirent pas encore aux catholiques français la pleine liberté de l'exercice de leur religion, du moins ils firent cesser de suite, presque entièrement, la persécution.

Nous pûmes donc enfin respirer un air un peu plus libre ; nous faisions alors (je vous demande pardon de la trivialité de la comparaison, mais elle vient si bien à mon sujet que je n'ai pas le courage d'y renoncer)
% Source : PDF original, vue 020, côté gauche, page manuscrite 40.
nous faisions comme les souris ; lorsque le gros chat qui leur faisait peur a disparu, elles sortent de leurs trous ; de même nous nous émancipions chaque jour davantage ; les prêtres ne se condamnaient déjà plus à une retraite aussi sévère ; nous nous réunissions en plus grand nombre dans les oratoires ; à l'occasion d'une fête solennelle ou de quelques cérémonies particulières, d'une première communion par exemple, on osait décorer les autels, brûler un peu d'encens, chanter à mi-voix les louanges du Seigneur sans craindre que l'odeur ou le bruit n'éveillât l'attention des persécuteurs. Bientôt on fit un pas de plus dans cette carrière de liberté : jusqu'ici tous les oratoires étaient dans des maisons particulières ; pour y être admis il fallait en conséquence obtenir l'agrément des personnes chez lesquelles ils se trouvaient ; on en ouvrit de publics, beaucoup plus vastes et on commença à y faire les offices aussi régulièrement que dans une église de paroisse.

Nous l'avons déjà remarqué, jusqu'à présent les jeunes gens fidèles à la religion avaient été isolés les uns des autres. Oh ! comme ils s'estimèrent heureux de pouvoir enfin apprendre à se connaître en se rencontrant aux pieds des autels ; il vous serait maintenant difficile, Messieurs, de pouvoir comprendre ce qui se passait réciproquement dans nos âmes, lorsque nous nous apercevions pour la première fois dans la maison du Seigneur. Pour s'en faire une idée, il faudrait avoir
% Source : PDF original, vue 020, côté droit, page manuscrite 41.
éprouvé comme nous tout ce qu'il y avait d'amer dans l'isolement auquel nous avions été condamnés. Je voudrais avoir des expressions assez vives pour faire passer dans vos cœurs les sentimens qui remplissaient les nôtres, lorsque au moment de la communion, plusieurs d'entre nous se levaient à la fois pour aller s'asseoir à la table sainte. Comme la piété des uns ravissait les autres ! C'était bien véritablement et bien entre nous la communion des saints, car le cœur de notre bon maître était là ; c'était lui qui nous unissait tous dans son amour. Aussi ces liaisons cimentées par le sang même du Sauveur devenaient inébranlables. Les cœurs s'étaient devinés ; les âmes s'étaient unies ; les corps se rapprochaient ; on sortait ensemble ; les mains se serraient ; une confiance réciproque faisait dilater les sentimens dont nous étions embrasés ; par de délicieux entretiens on se plaisait à les répandre au dehors, à les partager avec des nouveaux amis qu'on considérait déjà comme des frères. Dans ces conversations amicales on s'occupait surtout de la délivrance prochaine et du triomphe de la foi ; c'était une chaîne de charité que le Seigneur dans sa miséricorde formait à notre insu ; elle descendait du trône de son amour ; nos cœurs en étaient les anneaux ; elle devait nous soutenir jusqu'à la fin de nos jours dans le chemin glissant par lequel on parvient au vrai bonheur. Notre Sauveur nous façonnait ainsi
% Source : PDF original, vue 021, côté gauche, page manuscrite 42.
d'avance à devenir de véritables congréganistes ; ces douceurs de l'amitié chrétienne dont il remplissait nos âmes étaient comme un avant-goût de cette nouvelle terre promise dont sa bonté se préparait à nous mettre en possession, et qui n'est encore elle-même, malgré tous les ruisseaux de lait et de miel que nous y avons rencontrés, qu'une figure et une ombre de cette véritable terre promise pour laquelle nous soupirons et où nous parviendrons, je l'espère, après avoir traversé le désert de cette vie.

\section*{Première rencontre des amis chrétiens}
% Mention marginale de la page manuscrite 42.
C'est alors et de cette manière que j'eus le bonheur de vous rencontrer, bon et vénérable abbé Ripoud, lorsque le Seigneur vous a fait connaître pour nous donner un parfait modèle de mansuétude évangélique. C'est là comme l'aurore de cette douce liaison que plus tard la Ste Vierge, notre mère commune, a consacrée en nous adoptant ensemble pour ses enfants. La mort qui vous a mis en possession de votre récompense n'a certainement pas rompu notre union. Lorsque les jours de la tribulation seront achevés, j'ose l'espérer en m'appuyant sur la miséricorde infinie de mon Dieu, l'ami que vous avez laissé sur cette misérable terre viendra vous rejoindre dans les demeures éternelles. Là, mieux encore qu'ici-bas, comme vous m'en donniez si bien l'exemple, nous serons toujours dévoués à la Mère de mon Sauveur ; là nous louerons, nous bénirons, nous glorifierons notre Dieu pendant l'éternité ; nous verrons à découvert les motifs secrets de sa conduite à notre égard ; nous comprendrons
% Source : PDF original, vue 021, côté droit, page manuscrite 43.
alors parfaitement toute la grandeur de la grâce que nous avons reçue, lorsqu'une Congrégation qui devait dans ce monde remplir nos cœurs de tant de consolation et assurer notre salut dans l'autre a été formée pour nous.

Pardon, Messieurs, je m'oublie, je laisse mon âme s'épancher et ma plume trace en lignes sous sa dictée ; je me complais dans mes souvenirs ; mais il faut cependant continuer à suivre le cours des événements.

Lorsque Bonaparte eut pris les mesures les plus urgentes et les plus nécessaires pour affermir son pouvoir naissant, il se hâta de partir pour l'Italie afin de se remettre à la tête de son armée. Lorsque l'intérêt de l'Église l'avait exigé, le Seigneur, ainsi que nous l'avons vu, avait chassé cette armée de presque toute l'Italie ; mais aujourd'hui les décrets du Très-Haut sont accomplis, Rome est délivrée et le pape est intronisé ; la Providence veut maintenant augmenter la puissance de Bonaparte afin qu'il lui devienne possible d'exécuter les desseins qu'elle a sur lui ; aussi, après une bataille dont l'issue demeure incertaine pendant plusieurs jours, Dieu lui fit remporter à Marengo une victoire complète sur les Autrichiens. Ce nouveau succès eut pour résultat la conquête rapide du Piémont, de l'État de Gênes et de la Lombardie ; mais les États du pape furent pour lors respectés.

Bonaparte réunit le Piémont et l'État de Gênes à la France ; il recomposa la République cisalpine avec les États enlevés à la maison d'Autriche et convoqua à Lyon, sous le nom de Consulte cisalpine, les notabilités de cette nouvelle
% Source : PDF original, vue 022, côté gauche, page manuscrite 44.
République pour s'en faire nommer président.

\section*{Signature du concordat}
% Mention marginale de la page manuscrite 44.
Lorsque Bonaparte crut avoir ainsi consolidé son pouvoir, il prit la résolution de rétablir la religion en France et il ouvrit à cet effet des négociations avec le St Siège. Le St Père s'empressa d'envoyer à Paris, pour suivre une affaire si intéressante pour la paix de l'Église, deux hommes d'un mérite distingué, Mgr Spina, alors archevêque \textit{in partibus}\footnote{Traduction : « Dans les contrées. » L’expression accompagne ici le titre épiscopal. (Note du transcripteur).} de Corinthe et depuis archevêque de Gênes, et le P. Caselli, savant religieux qui plus tard a été cardinal.

Maintenant, Messieurs, permettez-moi de suspendre un instant mon récit, car j’ai besoin de me livrer avec vous à quelques réflexions. Il me semble que nous ne devons perdre aucune occasion de ranimer notre foi en cherchant à reconnaître dans les événements qui se passent sous nos yeux la trace évidente de l'action de la Providence. Or ici Dieu me paraît avoir voulu lui-même soulever un coin du rideau qui dérobe ordinairement à nos yeux les ineffables secrets, afin que nous puissions appercevoir à découvert les ressorts admirables par lesquels sa suprême sagesse imprime le mouvement aux affaires de ce monde. Nous ne chercherons pas à pénétrer le motif qui a fait agir Bonaparte dans cette circonstance ; c'est une cause seconde qui ne vaut pas la peine que nous nous y arrêtions ; mais nous porterons un œil attentif sur les diverses circonstances qui ont préparé le grand événement dont nous nous occupons et sans lesquelles il n'aurait pas pu avoir lieu. Pour parvenir donc à vaincre les obstacles qui pouvaient s'opposer au rétablissement de la religion en France,
% Source : PDF original, vue 022, côté droit, page manuscrite 45.
il a fallu d'abord que l'armée révolutionnaire fût expulsée d'Italie afin que le conclave pût s'assembler et donner un chef à l'Église ; il a fallu qu'avant la délivrance de l'Italie et la défaite de l'armée française, Bonaparte eût été rappelé, car pour nous servir d'une expression aujourd'hui consacrée par l'usage, s'il se fût trouvé à la tête de l'armée vaincue il fût devenu un homme impossible ; il a fallu ensuite qu'au moment où Bonaparte, échappé de l'Égypte et revenu en France comme un fugitif, il y fût accueilli comme un triomphateur ; il a fallu encore qu'au 18 brumaire, lorsque Bonaparte, au Conseil des Cinq-Cents, entendant les cris de hors la loi proférés contre lui, pâlissait, perdait contenance et était prêt à abandonner la partie, il a fallu, dis-je, que les grenadiers qui l'accompagnaient relevassent son courage abattu en criant à leur tour aux députés : Hors la salle, tout en les faisant sauter par les croisées et assurassent ainsi le succès de la journée ; il a fallu enfin que la bataille de Marengo vînt donner au pouvoir nouvellement établi du premier consul la sanction de la victoire et lui permettre de soumettre tout ce qui l'entourait à sa volonté de fer, car il n'était pas facile de comprimer tous ces impies qui ne pouvaient voir qu'en frémissant de rage une religion qu'ils croyaient avoir renversée pour jamais se relever triomphante, revêtue par le courage héroïque des martyrs qu'elle vient d'enfanter d'une gloire nouvelle. Observez bien que l'absence d'une seule de ces circonstances aurait suffi pour que notre
% Source : PDF original, vue 023, côté gauche, page manuscrite 46.
pauvre France restât longtemps encore privée de sa religion. Mais ce n'est pas tout ; si Bonaparte eût conçu alors la malheureuse pensée d'exiger du pape comme condition \textit{sine qua non}\footnote{Traduction : « Sans laquelle non. » C’est-à-dire une condition indispensable. (Note du transcripteur).} du traité ces concessions que son orgueil lui fit réclamer plus tard avec tant de persistance, il eût forcé le St Père de rompre les négociations. La main de Dieu peut-elle se laisser apercevoir avec plus de clarté ? Oh ! Messieurs, remercions donc cette suprême Providence qui dispose tout avec force, mais aussi avec douceur pour le plus grand avantage de ses élus, et écrions-nous dans un profond sentiment d'admiration : \textit{À Deo factum est istud...}\footnote{Traduction : « Cela a été fait par Dieu. » (Note du transcripteur).}

Mais voici encore une nouvelle occasion de nous exciter à la reconnaissance, et ce que je vais ajouter ici doit faire palpiter tous vos cœurs, vous tous qui vous honorez d'être les enfants de Marie. La piété de nos rois avait consacré notre France à cette reine du Ciel ; les peuples avaient uni leurs vœux à celui de leurs souverains ; la France tout entière prosternée aux pieds de Marie lui avait dit dans sa confiance toute filiale : Oh ! Mère de Dieu ! Soyez aussi à jamais notre mère ; régnez sur nous et sur nos enfants... Ce fut au jour de son glorieux triomphe que ce doux pacte de famille fut conclu avec la patronne de la France ; eh bien ! aujourd'hui Dieu veut que nous ne doutions pas que ce traité a été ratifié dans le ciel et que Marie nous a tous adoptés pour ses enfants, et pour nous démontrer jusqu'à l'évidence que c'est à l'intervention de notre bonne mère que nous devons le bonheur de conserver la foi, il permet que ce concordat
% Source : PDF original, vue 023, côté droit, page manuscrite 47.
qui doit nous rendre le libre exercice de notre religion soit signé aussi au jour du triomphe de Marie, 15 août 1801. Il me semble qu'ici les faits parlent tout seuls ; je plaindrais les cœurs froids qui ne comprendraient plus ce langage, mais ce ne seraient pas des cœurs congréganistes, des cœurs de la famille de Marie ; aussi vous le répéterez tous avec moi : Oui, Marie sera toujours notre mère. Oui, toujours nous remercierons Dieu de l'avoir glorifiée. Toujours nous nous confierons en sa puissante protection ; toujours nous nous réjouirons de lui appartenir ; toujours nous proclamerons le beau privilège de sa conception immaculée. Toujours enfin, en nous rappelant tous ses bienfaits, nous célébrerons dans les transports de la plus vive allégresse son glorieux triomphe. Vous vous souvenez sans doute que lorsque dans la folie de son orgueil Napoléon voulut essayer à l'ombre de son St patron d'usurper les honneurs qu'un peuple fidèle rendait à la mère de Dieu, ce projet insensé n'eut d'autre résultat que d'augmenter encore, s'il eût été possible, le concours des chrétiens qui se pressaient autour des autels de Marie.

Nous avancions à grands pas vers le rétablissement de la religion ; nous ne savions rien de positif sur l'issue des négociations, car tout se traitait dans le secret entre les deux puissances, et la presse, alors captive, ne livrait rien à la publicité. Néanmoins les heureuses nouvelles commençaient à transpirer et nous comblaient de joie ; nos cœurs s'épanouissaient et s'ouvraient à l'espérance. Essayons si par quelques comparaisons
% Source : PDF original, vue 024, côté gauche, page manuscrite 48.
je pourrais vous faire un peu comprendre ce qui se passait alors.

\section*{Espérances des catholiques}
% Mention marginale de la page manuscrite 48.
Figurez-vous un voyageur traversant une vaste forêt au milieu des ténèbres épaisses d'une nuit obscure ; il ne sait comment diriger ses pas incertains ; tout à coup une lueur encore presque insensible commence à lui faire espérer que le jour ne tardera pas à paraître. Peu à peu cette lueur augmente ; déjà les objets dont il est environné s'aperçoivent plus distinctement. Bientôt cette lumière naissante n'est plus une lumière douteuse ; les ombres de la nuit ont entièrement disparu ; enfin l'aurore brille à l'orient et annonce le retour prochain de l'astre bienfaisant qui va ranimer toute la nature. C'est ainsi que notre pauvre France tristement assise à l'ombre de la mort, sans culte et sans Dieu, pouvait enfin concevoir l'espérance fondée d'être bientôt éclairée de nouveau et ranimée par les rayons du soleil de justice.

Figurez-vous encore un homme surpris par un orage affreux au milieu d'une campagne déserte ; les nuages se sont amoncelés ; le ciel couvert d'épaisse nuées ne laisse plus parvenir jusqu'à lui un seul rayon de lumière ; un vent impétueux s'élève ; l'éclair sillonne la nue ; la foudre gronde à travers l'espace ; la grêle tombe avec fracas ; au milieu de cette effroyable tempête la nature entière semble s'abîmer. Le malheureux tremblant et consterné craint à chaque instant d'être écrasé
% Source : PDF original, vue 024, côté droit, page manuscrite 49.
par la foudre ; mais tout à coup le vent s'apaise ; le sifflement horrible de la grêle cesse ; les nuées s'entrouvrent, s'éclaircissent ; déjà le tonnerre ne gronde plus que dans le lointain ; le jour renaît ; les oiseaux recommencent leurs chants joyeux ; tout paraît reprendre une nouvelle vie ; l'infortuné respire. Si son cœur est encore attristé par l'aspect des nombreux débris qui attestent le passage de l'orage, l'espoir d'un temps meilleur le réjouit et le console ; s'il trouve alors sur son chemin un compagnon d'infortune qui comme lui vient d'échapper à la tempête, je vous laisse à penser ce qui se passe dans le cœur de ces deux hommes ; comme ils se hâtent de se féliciter, de se réunir et de se prêter un mutuel secours ; comme chacun d'eux se réjouit d'avoir trouvé un ami avec lequel il peut marcher sans crainte, tout en suivant le chemin qui doit le ramener dans la patrie ; c'est là que sous l'abri protecteur du toit paternel il sera pour toujours à l'abri de l'orage.

Tels étaient donc les sentimens dont nos cœurs étaient alors pénétrés ; tels furent aussi, je pense, ceux qu'éprouvèrent il y a quinze siècles les fidèles échappés à la cruelle persécution de Dioclétien ; comme nous, ils se livraient à une joie bien pure en voyant approcher le moment des miséricordes du Seigneur. C'est ainsi que Dieu dans sa bonté, après avoir épuisé sa justice, fait succéder le calme à la tempête ; alors ceux qui sont demeurés fidèles surabondent de consolations et des mains secourables
% Source : PDF original, vue 025, côté gauche, page manuscrite 50.
se présentent de toute part pour aider à se relever ceux qui ont eu le malheur de se laisser abattre.

Ne soyez pas étonnés, Messieurs et amis, si je mets quelque chaleur à vous retracer le souvenir de ces grands événements. Ils m'ont si fortement impressionné que malgré les années nombreuses qui se sont écoulées depuis et les glaces de l'âge qui m'environnent aujourd'hui je ne puis pas encore y arrêter une pensée sans que les sentimens qu'ils ont excités dans mon âme ne se reproduisent dans toute leur vivacité.

Ces sentimens du reste étaient alors partagés par une grande partie de la population lyonnaise. Tous les jours on s'apercevait davantage combien, malgré tous les efforts de l'impiété trop longtemps victorieuse, Lyon était demeuré catholique. Entre autres faits qui pourraient venir ici à l'appui de mon assertion, je me bornerai à vous citer le spectacle imposant que présenta notre ville au jour du Jeudi saint 1801. Nous n'étions alors en possession d'aucune église, et la population entière était sur pied pour la pieuse cérémonie des stations, tout comme cela se pratique à présent ; seulement au lieu de diriger ses pas d'une église à l'autre on allait de maison en maison ; il suffisait à un catholique de suivre la foule pour arriver à des oratoires dont plusieurs jusqu'à ce jour avaient pu lui être inconnus. On montait ou on descendait les nombreuses marches d'un escalier, assez souvent jusqu'au quatrième étage, toujours entouré des
% Source : PDF original, vue 025, côté droit, page manuscrite 51.
fidèles qui allaient ou qui venaient adorer notre Seigneur à un reposoir élevé quelquefois dans un galetas par les soins d'un pauvre ouvrier ou d'une dévideuse. Pendant que les nombreux oratoires catholiques se remplissaient ainsi de minute en minute par une multitude sans cesse renouvelée de fervents chrétiens, les églises schismatiques présentaient vainement au public leurs grandes portes ouvertes ; elles n'en restaient pas moins désertes, car il ne se trouvait personne qui voulût y entrer. J'ai vu ce jour-là un peuple nombreux se presser à l'entrée de l'allée étroite d'une maison de la place St Nizier pour parvenir à un modeste reposoir qui se trouvait dans l'intérieur de cette maison. Cela formait un singulier contraste avec la solitude dans laquelle gémissait tout près de là la vaste église qui servait encore de cathédrale aux constitutionnels.

Tout se rétablissait peu à peu et rentrait dans l'ordre. Le grand séminaire fut vers ce temps-là réintégré à la Providence ; on commença aussi à établir dans le diocèse plusieurs petits séminaires pour préparer d'avance par de bonnes études et l'habitude d'une vie pieuse de bons sujets pour le sacerdoce.

Les vicaires généraux accueillirent avec joie l'offre qui leur fut faite par les frères de la doctrine chrétienne de choisir la ville de Lyon pour y réunir leur société dispersée par la Révolution. Le conseil diocésain comprit quels immenses avantages tout le diocèse pouvait retirer de la présence de ces saints religieux. Les frères furent installés à Lyon
% Source : PDF original, vue 026, côté gauche, page manuscrite 52.
vers la fin de 1801 ou au commencement de 1802. Cet établissement, plus tard par la protection du cardinal Fesch, a pris un grand accroissement, s'est étendu de là, comme vous le savez, dans toute la France et a forcé ceux-là même qui détestent ces modestes instituteurs parce qu'ils sont au nombre des soutiens de la religion à rendre justice à la supériorité de leurs talents pour l'éducation de la jeunesse.

\section*{Arrivée des Pères de la foi à Lyon}
% Mention marginale de la page manuscrite 52.
Lyon fit encore cette année une autre acquisition. Celle-là se lie d'une manière directe et immédiate à l'objet qui nous occupe essentiellement dans ce précis ; je veux parler de l'arrivée des P. P. de la foi dans notre ville. Vous savez tous sans doute qu'on a appelé de ce nom une société de jeunes ecclésiastiques formée hors de France pendant la Révolution avec le projet de faire revivre autant qu'il était en elle, en se conformant à la règle que lui avait donnée son fondateur St Ignace, l'ordre célèbre des jésuites. Dès que la tranquillité avait commencé à paraître vouloir se rétablir en France, les religieux français de cet ordre nouveau s'étaient rendus à Paris où ils avaient formé dans le secret un premier établissement ; ils songèrent à étendre leur ordre. Le P. Varin, leur provincial, jeta les yeux sur Lyon pour y établir une seconde maison. Tout pouvait lui faire espérer que dans cette ville des martyrs, chef-lieu d'un diocèse habilement gouverné, il trouverait toutes les ressources nécessaires pour parvenir à son but ; il fit donc un voyage à Lyon, n'y
% Source : PDF original, vue 026, côté droit, page manuscrite 53.
passa que peu de jours, s'entendit avec les vicaires généraux, et dès qu'il fut de retour à Paris il se hâta d'envoyer pour former la maison de Lyon le P. Roger, auquel il donna pour compagnon le P. Barat, homme rempli de science et dont la modestie égalait le savoir. Plusieurs de leurs confrères vinrent au commencement de 1802 se joindre à ces deux fondateurs, et ils commencèrent dès lors à vivre en communauté. Je n'ai pas besoin de vous dépeindre le P. Roger, Messieurs ; presque tous vous l'avez connu. Vous savez donc tout comme moi qu'à un zèle ardent il joignait une douceur évangélique avec laquelle il attirait à lui tous les cœurs, et qu'il savait merveilleusement se servir de cet ascendant pour les porter à la pratique de la vertu ; aussi il ne lui fut pas difficile d'acquérir promptement la confiance des Lyonnais ; en peu de temps il devint le conseil et l'ami d'un grand nombre de fidèles. Dieu le préparait ainsi à son insu, tout comme au nôtre, à devenir bientôt un père pour nous tous. Il lui était réservé d'être le centre autour duquel devaient venir se grouper tous ces chrétiens auxquels le Seigneur se proposait de confier bientôt la noble mission de concourir au soutien et à la défense de la foi en secondant, en véritables troupes légères, par l'élan de leur zèle et de leur charité, l'action régulière des pasteurs de l'Église.
